% ---
% id: meat-core/ass:hash
% kind: ass
% title: "Ideal two-universal privacy amplification"
% summary: >-
%   Privacy amplification uses an ideal two-universal family, collision
%   probability exactly 2^−ℓ for distinct inputs, and error verification a
%   two-universal family. The first is what leftover hashing reads, the
%   second what the verification collision bound reads.
% prominence: standard
% status: accepted
% uses: [def:postprocessing-model]
% assumes: []
% refs:
%   - "tupkary2026decoy, Lemma 6.5 and Rem. 6.1"
%   - "kamin2025renyi, Thm. 4"
%   - "tupkary2026decoy, Def. 3.5"
% provenance: newly-derived
% short_id: A-HASH
% verification: never
% category: post-processing
% affects: [secrecy, correctness]
% user_statement: >-
%   Privacy amplification uses a hash family in which any two different
%   inputs collide with probability exactly 2^−ℓ (an ideal two-universal
%   family, such as Toeplitz hashing), and error verification uses a
%   two-universal hash family. If your privacy-amplification family is only
%   two-universal (collision probability at most 2^−ℓ), shorten the key by
%   one further bit.
% lean:
% ---

The R\'enyi leftover-hashing lemma \cite[Thm.~4]{kamin2025renyi} is applied
under the reading of \cite[Lemma~6.5 and Rem.~6.1]{tupkary2026decoy}, which
asks for an ideal family; the correctness argument uses the collision
probability of the verification family only as an upper bound. The two
notions are those of \cite[Def.~3.5]{tupkary2026decoy}.

\begin{assumption}\label{ass:hash}
(i) The privacy-amplification family of
Definition~\ref{def:postprocessing-model}, applied to $Z_1^n$, is ideal
two-universal: for any two distinct inputs, the probability under the
family's law that their $\ell$-bit values coincide is exactly $2^{-\ell}$;
(ii) the error-verification family of
Definition~\ref{def:postprocessing-model} is two-universal: for any two
distinct inputs, the probability under the family's law that their
$\lceil\log_2\frac1{\eev}\rceil$-bit values coincide is at most
$2^{-\lceil\log_2(1/\eev)\rceil}$.
\end{assumption}
